
Despite LoRA's widespread adoption as the parameter-efficient fine-tuning method of choice, practitioners universally default to rank $r=16$ without systematic justification. This one-size-fits-all approach may leave performance unrealized at scale. As language models grow from billions to hundreds of billions of parameters, the question becomes pressing: should LoRA rank scale with model size?

Prior work on LoRA has focused on the method itself---low-rank decomposition of weight updates, $\alpha /r$ scaling factors, and which layers to adapt---rather than how optimal rank varies across model scales. This gap leaves practitioners without guidance: a 12B model may benefit from rank$=64$ while a 1B model saturates at rank$=16$, yet both typically receive the same default configuration.

This work addresses this gap with a systematic empirical study across the Pythia model family (1B to 12B parameters). The primary finding is that optimal LoRA rank scales sub-linearly with model size: $r_{\text{opt}} \propto N^{\alpha}$ where $\alpha < 1$. On SQuAD-v2, $\alpha = 0.82$ (95\% CI: [0.73, 0.92]). This sub-linear relationship suggests that task-relevant parameter subspaces grow more slowly than overall model capacity.

Beyond the scaling law itself, this study documents two additional findings:

\begin{enumerate}
\item \textbf{Phase transition in rank sensitivity}: Larger models (12B) exhibit approximately $2\times$ higher sensitivity to rank selection than smaller models (1B), meaning the penalty for suboptimal rank increases with scale.
\end{enumerate}

\begin{enumerate}
\item \textbf{Inverted attention entropy}: Contrary to initial predictions, larger models show lower attention entropy (more focused attention patterns), not higher.
\end{enumerate}

An important limitation is that the scaling exponent $\alpha$ is task-dependent. Single-hop QA (SQuAD-v2) yields $\alpha \approx 0.82$ while multi-hop reasoning (HotpotQA) yields $\alpha \approx 0.30$. There is no universal scaling law---task complexity modulates the relationship.

The contributions are:

\begin{itemize}
\item Empirical characterization of $r_{\text{opt}}$ vs model size $N$ under controlled conditions
\item Demonstration that rank sensitivity increases with scale
\item Documentation of task-dependency, preventing overclaiming of a universal law
\item Observation that attention entropy decreases with scale
\end{itemize}
